\begin{textsample}{POS Dim 2 – go – Score 29.00 – go/go\_inf\_28.txt}  \label{ex:f2_pos_173}
V.
EDUCAÇÃO GOIANA : TEMAS CONTEMPORÂNEOS E DIVERSIDADES Com o intuito de atingir o maior quantitativo e diversidade de estudantes e de instituições que fazem parte da educação em Goiás, sejam elas públicas ( \textbf{federais}, \textbf{estaduais} e \textbf{municipais} ) ou privadas, o Documento Curricular para Goiás Ampliado propõe uma prática inclusiva para legitimar o direito : “ a educação com \textbf{qualidade} para todos e todas ”, conforme garante o Art. 26º da Declaração Mundial dos Direitos Humanos ( 1948 ), o Art. 205º da Constituição Brasileira ( 1988 ) e o Art. 1º da Lei de \textbf{Diretrizes} e Bases da Educação Nacional ( 1996 ).
Nesta perspectiva, esse documento buscou abranger particularidades, singularidades e especificidades da educação em Goiás, desde a Educação Infantil ao Ensino Fundamental, contemplando os estudantes dos 246 \textbf{municípios} goianos, tanto das áreas urbanas como das áreas rurais, e de todas as faixas etárias.
A preocupação em relação à pluralidade na educação goiana está na compreensão de um conjunto de fatores que se interligam e interagem contribuindo assim, cada um da sua maneira, com estas particularidades.
Dentre estes fatores pode -se relacionar : a extensão territorial do estado ; a população \textbf{superior} a seis milhões de pessoas, sendo ela de origem indígena ou de migrantes de outras regiões do Brasil e de outros países, isto desde o período colonial ; o crescimento urbano ; o avanço do agronegócio e as fronteiras abertas para receber cada vez mais migrantes, incluindo, recentemente, os refugiados estrangeiros.
É válido ressaltar que toda essa pluralidade merece atenção dos professores e das escolas em seus planejamentos e nas suas práticas docentes.
Afinal, essa heterogeneidade está contemplada no DC-GO por meio das habilidades dos componentes curriculares e estão garantidas em \textbf{políticas} públicas e documentos oficiais vinculados à educação, conforme são explicitadas a seguir.
Educação de Tempo Integral A educação goiana se molda a partir das necessidades do estado e a criação das instituições de tempo integral vem como resposta aos interesses sociais e trabalhistas deste tempo.
O estado conta com 52 \textbf{municípios} goianos com 197 escolas de período integral, sendo elas públicas ou privadas.
Existem as escolas de período integral ( os estudantes permanecem na instituição nos turnos matutino e vespertino ), as escolas integrais ( os estudantes permanecem nos três turnos e podem residir na escola ), as escolas integrais com contraturno ( em um turno o estudante estuda e no outro desenvolve atividades esportivas, artísticas e de descanso ) e os berçários-escola integrais que realizam a Educação Infantil e podem ter dinâmicas de \textbf{horários} diferenciadas, buscando \textbf{atender} todas as especificidades desta etapa da educação.
A educação de tempo integral possibilita aos estudantes mais tempo para desenvolver as habilidades previstas no DC-GO, favorecendo ainda mais a formação de estudantes críticos e autônomos.
Educação Quilombola Rural e Urbana Em diversas regiões do Brasil, entre os séculos XVI e XVIII, surgiram os quilombos que foram criados como refúgios de pessoas que escapavam da repressão/submissão a elas impostas durante o período escravocrata brasileiro.
Inicialmente, esses quilombos eram constituídos em sua maioria por pessoas negras, descendentes de africanos que vieram forçados para o Brasil e aqui foram escravizadas.
Eles tinham a primordial função de esconderijo, porém, também se tornaram locais para as pessoas que neles se abrigavam desenvolverem suas práticas e manifestações culturais, possibilitando, assim, que essas manifestações fossem repassadas para outras gerações.
Com o decorrer do tempo, estes quilombos, que eram comunidades isoladas, passaram a se relacionar com pessoas pertencentes a outros grupos, proporcionando, de certa forma, uma ampliação da diversidade cultural.
Há, atualmente, trinta e três comunidades legalizadas e reconhecidas como Quilombos em Goiás, sendo três como Quilombos Urbanos e as demais Quilombos Rurais.
Sobre a educação escolar, é preciso compreender que existem as escolas quilombolas e estudantes quilombolas estudando em outras escolas espalhadas pelo território goiano.
No DC-GO o estudo sobre a especificidade do povo quilombola é garantido, visando a que os estudantes do estado possam conhecer, compreender e valorizar a história e a importância desse povo ; e reconhecê -los como brasileiros e goianos.
Educação Inclusiva O DC-GO reconhece a importância das Necessidades Educativas Especiais para a promoção de uma educação inclusiva real no estado, acredita na autonomia das escolas e dos professores, professores de apoio e intérpretes para observar cada realidade e aplicarem as metodologias e práticas pedagógicas também especiais, garantindo, assim, a aplicabilidade do \textbf{currículo}.
Necessidades Educativas Especiais ( NEE ) : intelectuais, sensoriais, psicológicas ( emocionais ), físicas e de acessibilidades merecem toda a atenção dos educadores de Goiás, que ao unir as \textbf{políticas} públicas de inclusão escolar, nacionais e \textbf{estaduais}, com o DC-GO, poderão alcançar a inclusão de fato.
Ao compreender que existem dificuldades de aprendizagem derivadas de fatores orgânicos e/ou ambientais, NEE permanentes ( exigem adaptações generalizadas do \textbf{currículo} escolar, devendo o mesmo ser adaptado às características do estudante, durante grande parte ou de todo o seu percurso escolar ) e NEE temporárias ( exigem modificações parciais do \textbf{currículo} escolar, adaptando -o às características do estudante num determinado momento do seu desenvolvimento ), cada instituição de ensino fará suas adaptações necessárias.
Ressalta -se, também, a importância do ensino de Libras e do Braile nas instituições, do diálogo claro acerca do tema no ambiente escolar, e das modificações arquitetônicas necessárias, por exemplo, a colocação de indicações/sinalizações, rampas e corrimãos, dentre outras adaptações, consolidando a ideia de um ambiente inclusivo que promova ainda mais a acolhida e a inclusão desses estudantes e demais pessoas da comunidade escolar.
Educação Escolar Indígena Como em todo o território brasileiro, o estado de Goiás também tem povos indígenas, escolas indígenas e muitos estudantes indígenas em situação de itinerância.
Para abarcar todos esses estudantes, em especial os de origem Avá Canoeiros, Tapuia, Xavante, Chiquitano e Jagañu, que são a maioria em Goiás, há professores intérpretes que podem ser das etnias ou não, e passam por formações especiais, a fim de manter as tradições e propiciar o processo educacional.
Para a FUNAI – Fundação Nacional do Índio ( 2009 ) os povos indígenas têm direito a uma educação escolar específica, diferenciada, intercultural, bilíngue/multilíngue e comunitária, conforme define a \textbf{legislação} nacional que fundamenta a Educação Escolar Indígena.
A \textbf{política} educacional guarda relações inerentes com outras \textbf{políticas} e ações, desenvolvidas pela FUNAI e por outros \textbf{órgãos} de governo, voltadas aos povos indígenas, como \textbf{políticas} voltadas à gestão territorial, à sustentabilidade, à saúde, etc.
Por isso, a harmonização dessas ações convergentes é fundamental para o estabelecimento de relações do estado com povos indígenas, reconhecendo e respeitando a autonomia destes e suas formas próprias de organização, inclusive as educacionais, com línguas orais, bem como diferenciadas metodologias de ensino e aprendizagem.
Nessa perspectiva, o DC-GO, em suas habilidades, reforça a necessidade do estudo sobre os povos indígenas do estado, visando a que nossos estudantes possam se reconhecer, enquanto indígenas, ou conhecer a importância deles na formação territorial, socioeconômica e cultural brasileira, em especial do estado de Goiás.
Educação de Crianças e \textbf{Adolescentes} com Distorção de Idade-Ano Para o MEC ( 1996 ), a distorção idade-ano é a proporção de alunos com mais de 2 anos de atraso escolar.
Diversos são os motivos desse atraso, cabe à instituição de ensino verificar e propor a melhor forma de resolver essa distorção, além de garantir o desenvolvimento das habilidades essenciais, demonstradas nesse documento, para cada ano escolar.
As principais causas da distorção em Goiás, apontadas em pesquisas, como as do IBGE ( 2010 ) são a evasão e o abandono escolar, todavia existem causas primárias que contribuem para estas, e apesar de muitas vezes estarem intimamente ligadas à situação socioeconômica do aluno ( trabalho infantil e \textbf{adolescente} ), isso nem sempre é fator determinante.
Uma das principais consequências da distorção idade-ano é o baixo desempenho dos estudantes em atraso escolar e a relação leitura/escrita/interpretação, quando comparados aos estudantes regulares e a \textbf{reprovação}, o que pode ser evidenciado pelos resultados inferiores aos esperados nas avaliações nacionais e \textbf{estaduais} do Ensino Fundamental.
Turmas especiais, avaliações específicas, \textbf{horários} diferenciados e estratégias próprias para esses estudantes são ações de extrema importância para a correção da distorção e para a \textbf{garantia} do aprendizado real.
O estado de Goiás possui \textbf{políticas} educacionais nesse sentido ( \textbf{estaduais} e \textbf{municipais} ) e incentiva todas as instituições a procurarem os \textbf{órgãos} responsáveis para realizarem essas correções, tão importantes para os estudantes goianos.
O DC-GO chama a atenção de professores e instituições de ensino para estarem sempre em acordo com o Estatuto da Criança e do Adolescente ( 1990 ).
Educação de Adultos e Idosos Desde 2003, através do Projeto Brasil Alfabetizado, o Ministério da Educação incentiva os estados brasileiros a promoverem a educação de adultos e idosos nas instituições de ensino a fim de garantir o Ensino Fundamental para esses estudantes, que trazem um \textbf{perfil} diferente das distorções, muitos nunca frequentaram uma escola ou pararam de estudar por muitos anos.
Essa realidade merece uma atenção diferenciada, com propostas, \textbf{horários} ( noturno e finais de semana ), metodologias, período de duração ( semestral e não anual ) e ações inerentes a esse público.
A educação de adultos e idosos busca a melhoria na \textbf{qualidade} de vida, aumento da autoestima e maior sociabilização desses estudantes.
As habilidades apresentadas no DC-GO podem ser alteradas e remodeladas para essas faixas etárias nos \textbf{currículos} escolares e planos diários de professores.
A arquitetura e mobiliário da escola também deve promover a inclusão de adultos e idosos para a promoção do processo de ensino e aprendizagem com mais significância.
Essa preocupação com a alfabetização de idosos perpassa o Estatuto dos Idosos ( 2003 ), que traz considerações importantes sobre a educação desse público, que merece toda a atenção dos professores e representam significativo número em Goiás.
Educação do Campo A educação de estudantes campesinos, em Goiás, conta com 76 escolas rurais e algumas extensões específicas e com temporalidades diferenciadas ( nem sempre são escolas fixas ).
Estão localizadas em áreas rurais ( fazendas, sítios e chácaras ), em acampamentos temporários e em assentamentos de terras vinculadas aos movimentos sociais de luta pela \textbf{reforma} agrária, como o MST – Movimento dos Sem Terra.
Com estruturas organizacionais específicas para esses estudantes, obedecendo os períodos de lavoura, colheita, pecuária, dinâmica climática e relações socioambientais ( ciclos agrícolas ), as habilidades do DC-GO devem ser repensadas para esse sujeito do campo, garantindo o \textbf{currículo}, vinculado ao seu cotidiano, suas práticas de vivência e seu ambiente.
Os professores dessas instituições precisam de formação e um olhar diferenciado aos estudantes campesinos, a fim de garantir a Educação Básica concreta, preservação dos seus modos de vida, os espaços e transportes para o ensino.
Educação Ambiental A Educação Ambiental está presente no DC-GO nas habilidades dos componentes de Geografia, História, Ciências da Natureza, Língua Inglesa e Matemática.
Socioambiental é um termo muito utilizado quando se pensa em Educação Ambiental, o termo engloba o homem na natureza.
O ambiente é onde o homem está e se relaciona socialmente, ou seja, em todos os espaços.
As crianças e os \textbf{adolescentes} precisam dessa visão integrada para se sentirem parte da natureza, assim o processo de educação ambiental acontecerá de forma mais real.
A Educação Ambiental Escolar deve promover mudanças de hábitos e de atitudes a partir de conhecimentos adquiridos.
Essas mudanças devem ser transformadas em ações mais corretas na escola, em casa e nos seus espaços de vivências e de lazer.
Uma forma de promover essas mudanças seriam aulas extraclasses, visitas \textbf{técnicas} e trabalhos de campo que auxiliam na visualização dos impactos socioambientais presentes nas áreas urbanas e rurais.
Educação no Trânsito A mudança de comportamento de condutores e de pedestres no trânsito pode e deve começar no ambiente escolar.
O Departamento de Educação para o Trânsito de Goiás ( DETRAN - 2018 ) realiza palestras para todas as faixas etárias.
São momentos que têm como objetivo fazer com que a criança e o jovem reflitam sobre suas próprias atitudes no trânsito.
Para a Educação Infantil e o Ensino Fundamental Anos Iniciais e Finais os temas são trabalhados de forma lúdica, fazendo com que a criança se identifique nas situações propostas.
Para tal é utilizada a contação de histórias, o uso de fantoches e a vivência na faixa de pedestre.
As crianças e os \textbf{adolescentes} são convidados a refletirem de forma mais crítica, pensando em sua realidade e no que podem fazer para modificá -la.
Essas ações, no entanto, não podem ser de responsabilidade somente do DETRAN, a Educação no Trânsito deve acontecer em todas as instituições de ensino utilizando todos os componentes curriculares, a fim de desenvolver nos estudantes, as habilidades que provoquem a mudança de hábitos, além do conhecimento acerca da gravidade da falta de educação no trânsito em nosso Estado.
Educação Fiscal e Financeira A Educação Fiscal é uma prática de cidadania que envolve o aprofundamento da relação estado e sociedade na fiscalização e na gestão dos recursos públicos.
O programa desenvolvido na Secretaria da Fazenda de Goiás ( SEFAZ ) tem o objetivo de levar à comunidade em geral os conteúdos referentes ao papel social dos tributos, a importância dos orçamentos para o bom funcionamento da \textbf{administração} pública, a alocação e ao controle dos recursos, entre outros temas relacionados à gestão e fiscalização das finanças públicas ( SEFAZ, 2018 ).
Além desse programa, é preciso que professores promovam a Educação Fiscal ( formação de estudantes que podem atuar como fiscais dos gastos públicos ) e a Educação Financeira ( formação econômica dos estudantes ).
Estas têm o intuito de desenvolver nos estudantes várias habilidades que se referem à economia, como a relação gastos, ganhos e a prática de reservas econômicas como poupanças e aplicações.
Esse tipo de educação é feita em várias escolas internacionais, mas no Brasil ainda é uma prática rara.
Educação Política e Eleitoral A Educação Política é um processo de transmissão de informações e de conhecimentos cuja finalidade é disponibilizar ao estudante um repertório que lhe permita compreender as nuances dos debates e de organização política no Brasil, em seu estado e \textbf{município}.
Possui também a função de capacitar crianças e \textbf{adolescentes} para participar ativamente da \textbf{política} e compreender o processo eleitoral brasileiro.
Politizar é uma habilidade extremamente importante a ser desenvolvida nos estudantes goianos para garantir a defesa de valores fundamentais à convivência democrática.
A \textbf{política} deve envolver tolerância às diferenças, direito à contradição, ética, responsabilidade e o reconhecimento do outro.
Essa educação, nas escolas, com crianças e \textbf{adolescentes} diminuem manifestações de ódio e discriminação com ataques físicos, orais e virtuais.
Educação para Grupos Juvenis – Tribos Urbanas Em muitas escolas goianas os \textbf{adolescentes} de Anos Finais do Ensino Fundamental se organizam em grupos, os grupos juvenis, também denominados de tribos urbanas.
Os jovens decidem se reunir grupalmente através de músicas, vestimentas, símbolos corporais ou gostos por jogos, esportes, danças, filmes e festas.
Em Mesquita e Maia ( 2007 ) há um levantamento dos grupos existentes em Goiás : hippies, colsplayers, grafiteiros, emos, gamers, nerds, geeks, roqueiros, metaleiros, punks, skatistas, jogadores de RPG, torcidas organizadas de futebol, pichadores, góticos, funkeiros, pagodeiros e sertanejos.
Nesse contexto, os autores sugerem que os professores busquem compreender melhor este comportamento grupal, suas ações coletivas e suas simbologias corporais e identitárias, pois fazem parte da transitoriedade da fase \textbf{adolescente} para a adulta não somente no Brasil, mas em diversas nações mundiais.
Trazer para o ensino os seus símbolos pode favorecer e facilitar o processo de desenvolvimento de habilidades propostas no DC-GO, bem como a melhoria nas relações sociais nas escolas, diminuindo até mesmo o bullying.
Educação Alimentar e Nutricional Educação Alimentar e Nutricional ( EAN ) é um campo de conhecimento e de prática contínua e permanente, transdisciplinar, \textbf{intersetorial} e multiprofissional que visa promover a prática autônoma e voluntária de hábitos alimentares saudáveis, no contexto da realização do Direito Humano à Alimentação Adequada e da \textbf{garantia} da Segurança Alimentar e Nutricional.
A prática da EAN ( que faz parte de um conjunto de estratégias criadas para promover a alimentação adequada e saudável dos estudantes ) deve fazer uso de abordagens e de recursos educacionais problematizadores e ativos que favoreçam o diálogo junto a indivíduos e grupos populacionais, considerando todas as fases do \textbf{curso} da vida, etapas do sistema alimentar e as interações e significados que compõem o comportamento alimentar, em especial na escola.
Os \textbf{currículos} da Educação Infantil e do Ensino Fundamental deverão incluir o assunto Educação Alimentar e Nutricional, como exige a LDB – Lei de \textbf{Diretrizes} e Bases na Educação Brasileira 13.666/2018.
A importância dessa lei é agir de modo a reduzir a obesidade infantil e \textbf{adolescente} além de assegurar informações sobre alimentação saudável aos cidadãos desde a infância.
O tema é de grande importância nos tempos atuais, em que adultos com pouca formação ou com hábitos alimentares inadequados terminam por reforçar o interesse de crianças e de \textbf{adolescentes} por uma dieta pouco nutritiva.
Educação em Comunidades de Migrantes Internacionais A educação em Goiás também se preocupa com as comunidades internacionais presentes no estado.
Nelas estão estudantes que necessitam de uma educação diferenciada, que respeite suas culturas e línguas maternas.
O DC-GO sugere que os \textbf{municípios} que possuem essas comunidades promovam essa educação singular para que os estudantes de origem internacional sintam -se incluídos no processo educacional e que suas práticas sejam repassadas aos demais estudantes e comunidade escolar.
A seguir apresenta -se uma relação de \textbf{municípios} goianos e a procedência de suas comunidades internacionais : Vianópolis/Estados Unidos, Rio Verde/Alemanha, Colônia de Uvá/Alemanha, Nova Veneza/Itália ; Goiânia/Japão e Anápolis, Síria, Líbano, Turquia.
Sexualidade e Cuidados com o Corpo A temática Sexualidade e Cuidados com o Corpo vem sendo discutida e permeada em diferentes componentes do DC-GO, especialmente nas Ciências da Natureza, embora possa ser abordada na Educação Infantil e Ensino Fundamental de maneira gradativa, cabendo ao professor mediar essa temática.
Apesar de ter uma importante função preventiva, a educação sexual não devia cumprir um papel meramente informativo, mas sim com foco no desenvolvimento do indivíduo no respeito por si próprio, e, consequentemente pelo outro.
É importante que o estudante se aproprie do conhecimento científico a respeito do corpo humano, sobre as condições de vida da população e sobre a importância de colocar em prática hábitos que contribuirão decisivamente para o cuidado de si próprio.
O estudante deve perceber que hábitos de higiene ( escovar os dentes, banho, entre outros cuidados com a limpeza corporal ) o ajudam a possuir melhor \textbf{qualidade} de vida, mantendo -se saudável.
Boas práticas de higiene promovem, sobretudo, melhor convivência, evitando desconforto e até mesmo baixo rendimento escolar.
Educação Prisional As unidades escolares que atendem a modalidade para a Educação Prisional em Goiás possuem organização específica para cada penitenciária, presídio ou espaços de prisões provisórias, se diferenciando no número de alunos para abertura de turmas, na matriz curricular e na \textbf{carga} \textbf{horária}, nas atribuições, funções e gênero dos professores, nos diários de classe, nos relatórios de aprendizagem, na formação de grupos de estudos, na classificação e reclassificação dos estudantes, nos períodos de matrícula, na duração do \textbf{semestre} letivo, na aprovação, na retenção e no aproveitamento de estudos.
A Educação Prisional cuida do processo educacional de detentos, dos filhos de detentos que estão com os pais em situação de cárcere e de presidiários em \textbf{regime} semiaberto ( estes podem ir a escolas próximas ao presídio de domicílio ).
Quanto aos processos pedagógicos para a orientação do trabalho dos professores, estes existem e acontecem nas unidades/extensões escolares.
O trabalho educacional em âmbito prisional está voltado para a construção e reconstrução de valores humanos, demandando a constituição de um ambiente inclusivo \textbf{culminando} na atuação coletiva e participativa dentro da escola/extensão escolar, com vias a \textbf{atender} às diferenças, sem sufocá -las.
Um ambiente inclusivo comprometido com valores é necessário na educação em prisões, visando atingir não somente os presidiários, mas os filhos deles e delas que vivem em cárcere.
Educação Hospitalar A Educação Hospitalar em Goiás atende aos aspectos legais do Ministério da Educação e do Ministério da Saúde e propicia mediação da aprendizagem em classe hospitalar aos estudantes da Educação Infantil e Ensino Fundamental que estejam impossibilitados temporariamente de \textbf{frequentar} a escola regular por motivo de tratamento médico ou convalescença.
Esse trabalho é desenvolvido pela Secretaria de Estado da Educação de Goiás desde agosto de 1999 e pelas secretarias \textbf{municipais}.
Existem em Goiás 43 classes hospitalares funcionando com Educação Hospitalar, com professores e salas de aulas especiais para os estudantes que merecem atenção especial e educação de \textbf{qualidade} onde estão em tratamento, praticando inclusive a escuta pedagógica.
No Brasil, a educação hospitalar é reconhecida por meio da criação de uma \textbf{legislação} para a criança e \textbf{adolescente} hospitalizado, através da resolução nº 41 de outubro de 1995, onde diz que a crianças e os \textbf{adolescentes} possuem o"direito de desfrutar de alguma forma de recreação, programas de educação para a saúde, acompanhamento do \textbf{currículo} escolar durante sua permanência hospitalar"( BRASIL, 1995 ).
A Lei de \textbf{Diretrizes} e Bases da Educação Nacional de 1996 também reforça esse atendimento educacional em hospitais.
O atendimento será feito em classes, escolas, ou serviços especializados sempre que, em função das condições específicas do aluno não for possível a sua integração nas classes comuns de ensino regular ( BRASIL, 1996 ).
A Secretaria de Educação Especial do MEC conceitua Classe Hospitalar como uma das modalidades de atendimento especial para crianças e \textbf{adolescentes} : (... ) ambiente Hospitalar que possibilita o atendimento educacional de crianças e jovens internados, que necessitam de educação especial ou que estejam em tratamento. ( BRASIL, 1994 ).
Em 2002 o Ministério da Educação, por meio da Secretaria de Educação Especial, elaborou um documento de estratégias e orientações para o atendimento nas classes hospitalares, assegurando uma Educação Básica.
Educação Hospitalar, segundo a Secretaria de Educação Especial, é o atendimento pedagógico educacional que ocorre em ambientes de tratamento de saúde, seja em internação, atendimento hospital-dia e hospital-semana ou serviços de atenção integral à saúde mental.
Educação para Refugiados A Constituição da República do Brasil e a Lei 9.474/97 funcionam como base legal para criação e implementação de \textbf{políticas} públicas que visam à assistência e à integração dos refugiados, independente de faixa etária e gênero.
Fazem parte deste processo de cidadania, assegurar a \textbf{efetivação} dos direitos econômicos, sociais e culturais, garantindo exclusivamente o direito ao trabalho, à saúde e à educação.
O gabinete de Assuntos Internacionais do Governo do Estado de Goiás, localizado em Goiânia, vem dando assistência institucional e consular para esses refugiados, através de sua \textbf{gerência} de atração de investimentos e assuntos consulares.
Acredita -se na importância de ressaltar que refugiados são pessoas que se encontram fora do seu país por causa de fundado temor de perseguição por motivos de raça, religião, nacionalidade, opinião política ou participação em grupos sociais, e que não possa ( ou não queira ) voltar para casa.
São também pessoas obrigadas a deixar seu país devido a conflitos armados, violência generalizada e violação massiva dos direitos humanos, como crise econômica.
No que tange à educação, todos os refugiados em idade escolar devem ter acesso ao sistema de educação pública, sendo regularmente matriculados.
O refugiado que não teve seus estudos concluídos no país de origem será orientado sobre a possibilidade de sua continuidade.
O refugiado receberá orientações sobre procedimentos para a revalidação de documentos escolares, que deverão ser facilitados.
Àqueles cuja língua de origem não seja o português, em cooperação com instituições locais, são ministrados \textbf{cursos} de língua portuguesa, informações acerca da cultura brasileira e noções básicas da região onde foi encaminhado. ( ESTATUTO DO REFUGIADO/ACNUR, 1951 ).
O estado de Goiás tem recebido atualmente refugiados sírios, haitianos e venezuelanos, mas traz um histórico de receber outras nacionalidades.
O DC-GO garante em várias habilidades, de diferentes componentes, o respeito e valorização destas pessoas no estado, além de, incentivar aos professores uma recepção especial para com estes estudantes de Educação Infantil e Ensino Fundamental, devem também ter um cuidado especial com as línguas maternas e o ensino da Língua Portuguesa para os que não possuem como primeira língua, o português.
O respeito ao refugiado dentro das instituições de ensino faz melhorar a \textbf{qualidade} de vida dos refugiados, aumentar a diversidade cultural entre estudantes e professores e diminui a xenofobia, possibilitando mais respeito ao estrangeiro, ao de fora.

% matched lemmas: administração, adolescente, atender, carga, culminar, currículo, curso, diretriz, efetivação, estadual, federal, frequentar, garantia, gerência, horário, intersetorial, legislação, municipal, município, perfil, política, qualidade, reforma, regime, reprovação, semestre, superior, técnico, órgão
\end{textsample}
